\documentclass[11pt,aspectratio=169]{beamer}

\usetheme{metropolis}
\usepackage[utf8]{inputenc}
\usepackage[T1]{fontenc}
\usepackage{booktabs}
\usepackage{tikz}
\usetikzlibrary{shapes.geometric, arrows.meta, positioning}
\usepackage{array}

\definecolor{mathblue}{RGB}{41,98,158}
\definecolor{nodeteal}{RGB}{15,110,86}
\definecolor{nodecoral}{RGB}{214,93,74}
\definecolor{nodepurple}{RGB}{130,90,160}
\definecolor{softgray}{RGB}{110,110,110}

\setbeamercolor{palette primary}{bg=mathblue, fg=white}
\setbeamercolor{frametitle}{bg=mathblue, fg=white}
\setbeamercolor{progress bar}{fg=nodecoral, bg=softgray!25}
\setbeamercolor{alerted text}{fg=nodecoral}
\metroset{sectionpage=progressbar, numbering=fraction}

\title{Structure d'un nœud de notre arbre}
\subtitle{Origine et justification de chaque champ}
\author{}
\date{}

\newcommand{\origine}[2]{%
  \begin{tikzpicture}[remember picture]
  \node[rectangle, rounded corners=4pt, fill=#1!15, draw=#1, line width=0.6pt, inner sep=6pt] {\small\textcolor{#1}{\textbf{Origine : #2}}};
  \end{tikzpicture}
}

\begin{document}

\maketitle

\begin{frame}{Sommaire}
  \tableofcontents
\end{frame}

\section{Légende}

\begin{frame}{Quatre origines possibles pour un champ}
  \begin{columns}[T,onlytextwidth]
    \begin{column}{0.24\textwidth}
      \begin{center}
      \colorbox{softgray!15}{\textcolor{softgray}{\textbf{\phantom{x}Neutre\phantom{x}}}}\\[4pt]
      \small structurel
      \end{center}
    \end{column}
    \begin{column}{0.24\textwidth}
      \begin{center}
      \colorbox{nodecoral!15}{\textcolor{nodecoral}{\textbf{\phantom{x}Détecteur\phantom{x}}}}\\[4pt]
      \small position + type grossier
      \end{center}
    \end{column}
    \begin{column}{0.24\textwidth}
      \begin{center}
      \colorbox{nodeteal!15}{\textcolor{nodeteal}{\textbf{\phantom{x}Géométrie\phantom{x}}}}\\[4pt]
      \small algorithme mathématique
      \end{center}
    \end{column}
    \begin{column}{0.24\textwidth}
      \begin{center}
      \colorbox{nodepurple!15}{\textcolor{nodepurple}{\textbf{\phantom{x}MLLM\phantom{x}}}}\\[4pt]
      \small sens précis
      \end{center}
    \end{column}
  \end{columns}
  \vspace{8mm}
  \begin{center}
  \small Chaque champ du nœud est rattaché à l'une de ces quatre origines --- jamais choisi au hasard.
  \end{center}
\end{frame}

\section{Les sept champs}

\begin{frame}{id}
  \origine{softgray}{neutre / structurel}
  \vspace{6mm}
  \begin{itemize}
    \item Identifiant unique du nœud
    \item Permet de faire correspondre deux arbres lors de l'évaluation (Tree Edit Distance, Container Match)
    \item Sans lui, impossible de relier un enfant à son parent sans ambiguïté
  \end{itemize}
\end{frame}

\begin{frame}{bbox --- position (x, y, largeur, hauteur)}
  \origine{nodecoral}{détecteur (feuille) / calcul (nœud interne)}
  \vspace{6mm}
  \begin{itemize}
    \item Pour une feuille : fournie directement par le détecteur
    \item Pour un nœud interne : \textbf{calculée} en agrégeant les bbox des enfants
    \item Le champ fondateur --- sans lui, ni découpage, ni rendu final
  \end{itemize}
\end{frame}

\begin{frame}{type grossier --- Texte / Image / Composant}
  \origine{nodecoral}{détecteur}
  \vspace{6mm}
  \begin{itemize}
    \item Fourni par le détecteur, sans raisonnement, uniquement sur les feuilles
    \item Sert d'\textbf{indice de départ} transmis au MLLM
    \item Le MLLM affinera ce type grossier en un sens précis, jamais l'inverse
  \end{itemize}
\end{frame}

\begin{frame}{parent}
  \origine{nodeteal}{algorithme géométrique}
  \vspace{6mm}
  \begin{itemize}
    \item Référence vers le nœud englobant
    \item Permet de remonter jusqu'à la racine de l'arbre
    \item Nécessaire pour fusionner deux groupes voisins quand un découpage s'avère accidentel
  \end{itemize}
\end{frame}

\begin{frame}{children --- liste ordonnée}
  \origine{nodeteal}{algorithme géométrique}
  \vspace{6mm}
  \begin{itemize}
    \item Définit littéralement la hiérarchie elle-même
    \item L'\textbf{ordre} porte le sens de lecture (gauche→droite, haut→bas)
    \item Indispensable pour générer un code dans le bon ordre visuel
  \end{itemize}
\end{frame}

\begin{frame}{type géométrique --- Feuille / Row / Column}
  \origine{nodeteal}{algorithme géométrique}
  \vspace{6mm}
  \begin{itemize}
    \item Indispensable à l'algorithme de découpage lui-même (XY-Cut)
    \item Précise si un nœud est terminal ou un conteneur, et dans quel sens il a été coupé
    \item Détermine ensuite le CSS à générer (flex-direction: row ou column)
  \end{itemize}
\end{frame}

\begin{frame}{label sémantique}
  \origine{nodepurple}{MLLM}
  \vspace{6mm}
  \begin{itemize}
    \item Description ou nom attribué par le MLLM, en tout dernier
    \item Rempli seulement une fois l'arbre géométrique entièrement figé
    \item Notre \textbf{contribution centrale} --- sans lui, l'arbre reste aussi muet que Prototype2Code
  \end{itemize}
\end{frame}

\section{Synthèse}

\begin{frame}{Vue d'ensemble des sept champs}
  \small
  \renewcommand{\arraystretch}{1.5}
  \begin{tabular}{@{}p{3.4cm}p{2.6cm}p{6.5cm}@{}}
    \toprule
    \textbf{Champ} & \textbf{Origine} & \textbf{Rôle} \\
    \midrule
    id & \textcolor{softgray}{Neutre} & Correspondance entre arbres \\
    bbox & \textcolor{nodecoral}{Détecteur} & Position exacte, agrégée pour les nœuds internes \\
    type grossier & \textcolor{nodecoral}{Détecteur} & Indice de départ pour le MLLM \\
    parent & \textcolor{nodeteal}{Géométrie} & Remontée et fusion \\
    children & \textcolor{nodeteal}{Géométrie} & Définit la hiérarchie et l'ordre \\
    type géométrique & \textcolor{nodeteal}{Géométrie} & Pilote le découpage et le CSS \\
    label sémantique & \textcolor{nodepurple}{MLLM} & Contribution centrale \\
    \bottomrule
  \end{tabular}
\end{frame}

\begin{frame}{Question ouverte avant d'augmenter la liste}
  \begin{center}
  \large Faut-il déjà ajouter des champs supplémentaires,\\ ou rester sur ces sept pour l'Option A ?
  \\[6mm]
  \small ex.\ une empreinte (hash) pour détecter les composants réutilisables --- pertinent seulement pour l'Option B
  \end{center}
\end{frame}

\begin{frame}[standout]
  Merci --- retours bienvenus
\end{frame}

\end{document}
